\chapter{Data type transformations}
\label{ch:type-transforms}

\begin{quote}
\emph{``Machine learning algorithms don't understand categories --- they understand numbers.''}
\end{quote}

% ============================================================
\section{Why transform data types?}

Most machine learning algorithms require numerical input. Raw datasets contain categorical variables (``Male'', ``Female''), ordinal variables (``Low'', ``Medium'', ``High''), and numerical variables on vastly different scales (age in years, income in thousands). Proper encoding and scaling are essential.

\begin{datatip}
Type transformation has two faces: \textbf{encoding} converts categories to numbers, and \textbf{scaling} puts numbers on comparable ranges. Both are required for most ML pipelines.
\end{datatip}

% ============================================================
\section{Categorical encoding}

\subsection{Label (ordinal) encoding}

\begin{minted}{python}
import pandas as pd
from sklearn.preprocessing import OrdinalEncoder

# Adult Census dataset
url = ("https://archive.ics.uci.edu/ml/machine-learning-databases/"
       "adult/adult.data")
cols = ["age", "workclass", "fnlwgt", "education", "education_num",
        "marital_status", "occupation", "relationship", "race",
        "sex", "capital_gain", "capital_loss", "hours_per_week",
        "native_country", "income"]
df = pd.read_csv(url, header=None, names=cols, na_values=" ?",
                 skipinitialspace=True)

# Ordinal encoding: for variables with a natural order
edu_order = [["Preschool", "1st-4th", "5th-6th", "7th-8th", "9th",
              "10th", "11th", "12th", "HS-grad", "Some-college",
              "Assoc-voc", "Assoc-acdm", "Bachelors", "Masters",
              "Prof-school", "Doctorate"]]

enc = OrdinalEncoder(categories=edu_order,
                     handle_unknown="use_encoded_value",
                     unknown_value=-1)
df["education_ord"] = enc.fit_transform(df[["education"]])
print(df[["education", "education_ord"]].drop_duplicates()
        .sort_values("education_ord"))
\end{minted}

\begin{warningbox}
Ordinal encoding imposes a numeric order. Using it on nominal categories (e.g., country names) implies ``France $<$ Germany $<$ Japan'', which is meaningless and can mislead models.
\end{warningbox}

\subsection{One-hot encoding}

\begin{minted}{python}
from sklearn.preprocessing import OneHotEncoder

# One-hot encoding: for nominal (unordered) categories
ohe = OneHotEncoder(sparse_output=False, drop="first",
                    handle_unknown="ignore")
encoded = ohe.fit_transform(df[["sex", "race"]])
encoded_df = pd.DataFrame(encoded,
                          columns=ohe.get_feature_names_out())
print(encoded_df.head())
print(f"New columns: {encoded_df.shape[1]}")
\end{minted}

\begin{minted}{python}
# pandas get_dummies: quick one-hot encoding
df_encoded = pd.get_dummies(df, columns=["sex", "race"],
                            drop_first=True, dtype=int)
print(df_encoded.columns.tolist()[-10:])
\end{minted}

\begin{preprocesstip}
Use \texttt{drop="first"} (or \texttt{drop\_first=True}) to avoid multicollinearity. With $k$ categories, you only need $k-1$ dummy variables --- the dropped category becomes the reference level.
\end{preprocesstip}

\subsection{Target encoding}

\begin{minted}{python}
# Target encoding: replace each category with the mean of the target
# Useful for high-cardinality columns (e.g., 100+ categories)
target_means = df.groupby("occupation")["income"].apply(
    lambda x: (x == ">50K").mean()
).to_dict()

df["occupation_target_enc"] = df["occupation"].map(target_means)
print(df[["occupation", "occupation_target_enc"]]
        .drop_duplicates().sort_values("occupation_target_enc"))
\end{minted}

\begin{warningbox}
Target encoding causes data leakage if applied before train/test split. Always fit on training data only, then transform both train and test. Use \texttt{sklearn.model\_selection.cross\_val\_predict} or K-fold target encoding to reduce overfitting.
\end{warningbox}

\subsection{Frequency and binary encoding}

\begin{minted}{python}
# Frequency encoding: replace category with its frequency
freq = df["native_country"].value_counts(normalize=True)
df["country_freq"] = df["native_country"].map(freq)
print(df[["native_country", "country_freq"]].head(10))

# Binary encoding: fewer columns than one-hot for high cardinality
# category-encoders library provides this
# pip install category-encoders
# import category_encoders as ce
# encoder = ce.BinaryEncoder(cols=["native_country"])
# df_binary = encoder.fit_transform(df)
\end{minted}

% ============================================================
\section{Numerical scaling}

\subsection{StandardScaler (z-score normalization)}

\begin{minted}{python}
from sklearn.preprocessing import StandardScaler

scaler = StandardScaler()
df[["age_scaled", "hours_scaled"]] = scaler.fit_transform(
    df[["age", "hours_per_week"]]
)

print(df[["age", "age_scaled", "hours_per_week", "hours_scaled"]].describe())
\end{minted}

\subsection{MinMaxScaler}

\begin{minted}{python}
from sklearn.preprocessing import MinMaxScaler

minmax = MinMaxScaler(feature_range=(0, 1))
df[["age_minmax"]] = minmax.fit_transform(df[["age"]])

# After scaling: min=0, max=1
print(df["age_minmax"].describe())
\end{minted}

\subsection{RobustScaler}

\begin{minted}{python}
from sklearn.preprocessing import RobustScaler

# Uses median and IQR: robust to outliers
robust = RobustScaler()
df[["capital_gain_robust"]] = robust.fit_transform(df[["capital_gain"]])
print(df["capital_gain_robust"].describe())
\end{minted}

\begin{preprocesstip}
Choose your scaler based on the data and the model:
\begin{itemize}
  \item \textbf{StandardScaler}: default choice for most algorithms (SVM, logistic regression, neural networks).
  \item \textbf{MinMaxScaler}: when you need bounded values (neural networks with sigmoid output).
  \item \textbf{RobustScaler}: when your data has significant outliers.
\end{itemize}
\end{preprocesstip}

% ============================================================
\section{Comparing scalers visually}

\begin{minted}{python}
import matplotlib.pyplot as plt
import numpy as np

fig, axes = plt.subplots(1, 4, figsize=(16, 4))

# Original
df["capital_gain"].hist(bins=50, ax=axes[0])
axes[0].set_title("Original")

# Standard
axes[1].hist(StandardScaler().fit_transform(df[["capital_gain"]]),
             bins=50)
axes[1].set_title("StandardScaler")

# MinMax
axes[2].hist(MinMaxScaler().fit_transform(df[["capital_gain"]]),
             bins=50)
axes[2].set_title("MinMaxScaler")

# Robust
axes[3].hist(RobustScaler().fit_transform(df[["capital_gain"]]),
             bins=50)
axes[3].set_title("RobustScaler")

plt.tight_layout()
plt.savefig("scaler_comparison.png", dpi=150)
plt.show()
\end{minted}

% ============================================================
\section{Discretization (binning)}

\begin{minted}{python}
# Equal-width binning
df["age_bin_width"] = pd.cut(df["age"], bins=5,
                              labels=["very_young", "young", "mid",
                                      "senior", "elderly"])

# Equal-frequency (quantile) binning
df["age_bin_quantile"] = pd.qcut(df["age"], q=5,
                                  labels=["Q1", "Q2", "Q3", "Q4", "Q5"])

# Custom bins from domain knowledge
df["age_group"] = pd.cut(df["age"],
                          bins=[0, 25, 35, 50, 65, 100],
                          labels=["<25", "25-34", "35-49", "50-64", "65+"])

print(df["age_group"].value_counts().sort_index())
\end{minted}

\begin{minted}{python}
# KBins discretization from sklearn
from sklearn.preprocessing import KBinsDiscretizer

kbd = KBinsDiscretizer(n_bins=5, encode="ordinal", strategy="quantile")
df[["hours_binned"]] = kbd.fit_transform(df[["hours_per_week"]])
print(df[["hours_per_week", "hours_binned"]].head(10))
\end{minted}

% ============================================================
\section{Power transforms}

\begin{minted}{python}
from sklearn.preprocessing import PowerTransformer

# Yeo-Johnson: works with positive and negative values
pt = PowerTransformer(method="yeo-johnson")
df[["capital_gain_yj"]] = pt.fit_transform(df[["capital_gain"]])

# Box-Cox: only works with strictly positive values
# pt_bc = PowerTransformer(method="box-cox")

# Log transform: manual alternative for skewed data
df["capital_gain_log"] = np.log1p(df["capital_gain"])

print(f"Original skew:    {df['capital_gain'].skew():.2f}")
print(f"Yeo-Johnson skew: {df['capital_gain_yj'].skew():.2f}")
print(f"Log skew:         {df['capital_gain_log'].skew():.2f}")
\end{minted}

% ============================================================
\section{Exercises}

\begin{exercise}
\begin{enumerate}
  \item Load the Titanic dataset. Apply one-hot encoding to \texttt{Sex} and \texttt{Embarked}. Apply ordinal encoding to \texttt{Pclass} (1st $>$ 2nd $>$ 3rd). Show the resulting DataFrame.
  \item For the Adult Census dataset, compare the effect of StandardScaler, MinMaxScaler, and RobustScaler on \texttt{capital\_gain}. Plot all three distributions.
  \item Implement target encoding for \texttt{native\_country} using 5-fold cross-validation to prevent leakage. Compare the encoded values with naive target encoding.
  \item Create age groups for the Titanic dataset using domain-meaningful bins (child, teenager, young adult, middle-aged, senior). Compute survival rate by age group.
  \item Apply a Yeo-Johnson power transform to all numerical columns in the Adult dataset. Compare skewness before and after.
\end{enumerate}
\end{exercise}

% ============================================================
\section{Chapter summary}

\begin{itemize}
  \item Categorical encoding converts non-numeric data to numbers: use ordinal encoding for ordered categories, one-hot for nominal, and target encoding for high-cardinality features.
  \item Scaling puts numerical features on comparable ranges: StandardScaler for general use, MinMaxScaler for bounded ranges, RobustScaler when outliers are present.
  \item Discretization (binning) converts continuous variables into categories, useful for non-linear relationships and interpretability.
  \item Power transforms (Yeo-Johnson, Box-Cox, log) reduce skewness and can improve model performance.
  \item Always fit transformers on training data only, then apply to test data.
\end{itemize}
